\documentclass[11pt]{article}
\usepackage[margin=2.2cm]{geometry}
\usepackage{amsmath,amssymb}
\usepackage{graphicx}
\usepackage{enumitem}
\usepackage{parskip}
\usepackage{booktabs}
\usepackage[hidelinks]{hyperref}

\graphicspath{{figures/}{../paper/figures/}}
\setlist{nosep, leftmargin=1.4em}

\begin{document}

{\large\textbf{Expiry-day microstructure and the settlement VWAP: progress notes}}\\
Follow-ups from the meeting of 21 Oct 2026\\
\small{(working notes; the full report has the details and all tests)}
\normalsize

\hrule
\vspace{0.8em}

\section*{Setup (unchanged)}
\begin{itemize}
  \item NSE order and trade files, cash market, 2022. Full year (242 sessions) for trades;
  order book rebuilt for the 12 monthly expiries + the day before each.
  \item 150 stocks: 50 liquid F\&O, 50 illiquid F\&O, 50 placebo (no derivatives, matched on
  turnover and price).
  \item New: F\&O order and trade files for Mar to Jun (4 expiries), used for the futures.
  \item Comparisons: expiry day vs.\ day before, same stock; regressions with stock and
  session fixed effects.
\end{itemize}

\begin{center}
\small
\begin{tabular}{llll}
\toprule
 & Hypothesis & Primary measure & Result \\
\midrule
BH1 & Settlement price is pushed & $|$settle VWAP $-$ mid at 15:00$|$ & not rejected ($p=0.21$) \\
BH2 & Window costlier to trade in & quoted spread & rejected ($p<0.001$) \\
BH3 & One-sided pressure & $|$OIB$|$ & not rejected ($p=0.54$) \\
BH4 & Order behaviour changes & orders cancelled within 1\,s & rejected ($p<0.001$) \\
\bottomrule
\end{tabular}
\end{center}

\section*{1. Reasons behind the results}
\begin{itemize}
  \item \textbf{BH1 (no push in stocks).} Stock F\&O settle by delivery, so moving the VWAP gains
  nothing: in the last minutes the expiring future tracks the live cash price (slope 0.91), not
  the VWAP. No pull towards option strikes. The index settles in cash, and there we do see it:
  Nifty 50 falls 6.2 bps vs.\ Nifty Next 50 over the window on index expiries ($p=0.008$).
  \item \textbf{BH2 (wider spread).} Driven by volatility, see section 3.
  \item \textbf{BH3 (no one-sided flow).} Flow at the close is two-sided (institutions closing
  and benchmarking; custodian share of window volume 59\% $\to$ 65\%), and futures positions are
  mostly rolled before expiry day (section 4).
  \item \textbf{BH4 (more fleeting orders).} Also mostly volatility: once we control for the rise
  in volatility, the increase is small and not significant.
\end{itemize}

\section*{2. Figure 2: is the next-morning fall news?}
\begin{center}
\includegraphics[width=0.95\textwidth]{window_path.pdf}
\end{center}
Checks:
\begin{itemize}
  \item drop stock-days with signs of news next morning (volume $>2\times$ median, or a very
  large opening gap);
  \item adjust for market news (stock beta $\times$ market's overnight move);
  \item drop 24 Feb 2022 (day of the Ukraine invasion: crash, then rebound next morning);
  \item estimate expiry by expiry.
\end{itemize}

\begin{center}
\small
\begin{tabular}{lcc}
\toprule
Extra reversal on monthly expiries & F\&O stocks & Placebo \\
\midrule
Baseline & $-0.26$ ($p=0.04$) & $-0.23$ ($p=0.007$) \\
Without news days & $-0.22$ ($p=0.08$) & $-0.09$ ($p=0.06$) \\
Beta-adjusted & $-0.15$ ($p=0.20$) & $-0.15$ ($p=0.006$) \\
Beta-adjusted, without news days & $-0.13$ ($p=0.34$) & $-0.06$ ($p=0.17$) \\
\bottomrule
\end{tabular}
\end{center}

\begin{center}
\includegraphics[width=\textwidth]{overnight_by_expiry.pdf}\\
\small Each monthly expiry (orange) vs.\ ordinary sessions (grey).
\end{center}

\textbf{Takeaway:} mostly information arriving after the close. It comes from a few expiries,
is just as strong in placebo stocks (nothing settles there), and is small and insignificant
after the two adjustments. One caveat: the morning-to-15:00 move is reversed much less than the
window move, so maybe not all of it is news. Checking against actual company announcements
would settle it.

\section*{3. BH2: spread and depth plots}
\begin{center}
\includegraphics[width=\textwidth]{bh2_spread_day.pdf}\\
\includegraphics[width=\textwidth]{bh2_depth_day.pdf}\\
\small Median across stocks, expiry day vs.\ day before (log scale). Window shaded.
\end{center}

\begin{center}
\includegraphics[width=0.9\textwidth]{bh2_window_zoom.pdf}\\
\small Last hour, minute by minute.
\end{center}

\begin{center}
\small
\begin{tabular}{lrrr}
\toprule
Window, expiry vs.\ day before & Spread & Depth & Volatility \\
\midrule
Liquid F\&O & $+9.8$\% & $-7.2$\% & $+45$\% \\
Illiquid F\&O & $+10.3$\% & $-4.8$\% & $+46$\% \\
Placebo & $+1.2$\% (n.s.) & $-0.5$\% (n.s.) & $+7$\% (n.s.) \\
\bottomrule
\end{tabular}
\end{center}

\textbf{Why:} the window is far more volatile on expiry days, only in F\&O stocks. Regressing each
stock's change in spread on its change in volatility: volatility explains about 70\% of the
widening and all of the fall in depth. Market makers quote wider and thinner when prices move
more.

\section*{4. Futures}
Rebuilt the order book of the expiring stock futures from the F\&O order files (4 expiries,
Mar to Jun, and the day before; 100 stocks, 84 on 28 Apr because one file part is an incomplete
download).

\begin{center}
\includegraphics[width=\textwidth]{futures_book_day.pdf}\\
\small Cash (left) vs.\ expiring future (right).
\end{center}

\begin{center}
\small
\begin{tabular}{lrrr}
\toprule
Window, expiry vs.\ day before & Spread & Depth & Volatility \\
\midrule
Cash (same stocks, Mar to Jun) & $+2.5$\% & $-1.0$\% & $+7.5$\% \\
Expiring future & $+31.9$\% & $-2.7$\% & $+39.9$\% \\
\bottomrule
\end{tabular}
\end{center}

\begin{itemize}
  \item The futures spread jumps in the window on expiry day. Volatility explains only about a
  fifth of it.
  \item Likely reason: it is the contract's last half hour. Most positions have already rolled
  (next month has 53\% of futures volume the day before, 59\% on expiry day), and anything left
  open turns into a delivery obligation, so market makers pull out.
  \item Still no one-sided pressure: futures $|$OIB$|$ in the window unchanged (0.18 vs.\ 0.18).
\end{itemize}

\section*{Next}
\begin{itemize}
  \item Company announcements data to test the news explanation directly.
  \item Futures book for more sessions, not only expiry and the day before.
  \item All of the above is in the updated full report.
\end{itemize}

\end{document}
